\section{Teorias de primeira ordem}\label{sec:teoriasPrimeiraOrdem}
\index{teoria da prova}
\index{sistema de Hilbert}
Até aqui, discutimos o que é uma linguagem de primeira ordem e aprendemos a identificar termos, sentenças e fórmulas, bem como tautologias.
O próximo passo é entender teorias, teoremas e demonstrações, ou seja, teoria da prova básica.

Para se provar teoremas, é necessário estabelecer axiomas (fórmulas que serão consideradas verdadeiras) e regras de inferência (maneiras de obter novas fórmulas a partir de fórmulas já estabelecidas como verdadeiras).

Os axiomas lógicos adotados variam muito de texto para texto.
Há quem diga que há tantas formas de axiomatizar a lógica quanto há textos de lógica.
Acreditamos que, em algum sentido, este texto somará 1 a essa contagem.
De qualquer modo, a grande maioria das formas é, em algum sentido, equivalente no caso da lógica de primeira ordem.
Não comentaremos outras formas de fazê-lo com profundidade.


Agora estamos prontos para enunciar os axiomas lógicos adotados neste texto.
Adotaremos um sistema de axiomas estilo Hilbert.
Ressaltamos novamente que esse não é o único modo de formalizar a dedução lógica.
Estamos adotando esse sistema por brevidade e simplicidade, uma vez que é um sistema de apresentação mais rápida e concisa.

\begin{metadefinition}[Axiomas Lógicos]\label{mdef:axiomasLogicos}\index{axioma!lógico}
    Os axiomas lógicos adotados neste texto são fórmulas de uma das seguintes formas.
    Abaixo, $A$, $B$ e $C$ são fórmulas, e $x, y$ são variáveis.
    \begin{enumerate}
        \item $A\rightarrow (B \rightarrow A)$.
        \item $(A \rightarrow (B \rightarrow C)) \rightarrow ((A \rightarrow B) \rightarrow (A \rightarrow C))$.
        \item $A \rightarrow (B \rightarrow (A \wedge B))$.\label{ax:conjuncao1}
        \item $(A \wedge B) \rightarrow A$.\label{ax:conjuncao2}
        \item $(A \wedge B) \rightarrow B$.\label{ax:conjuncao3}
        \item $A \rightarrow (A \vee B)$.
        \item $B \rightarrow (A \vee B)$.
        \item $(A \rightarrow C) \rightarrow ((B \rightarrow C) \rightarrow ((A \vee B) \rightarrow C))$.
        \item $\neg \neg A \rightarrow A$.
        \item $(A\leftarrow B) \rightarrow (A\rightarrow B)$.
        \item $(A\rightarrow B) \rightarrow (A\leftarrow B)$.
        \item $(A\rightarrow B) \rightarrow ((B\rightarrow A) \rightarrow (A \leftrightarrow B))$.
    \end{enumerate}
\end{metadefinition}

Ressalta-se que a Metadefinição~\ref{mdef:axiomasLogicos} nos dá um esquema de axiomas.
Esse esquema não consiste em um número finito de axiomas, e sim um número finito de esquemas.

A seguir, definiremos a noção de teoria, de demonstração, e de teorema.

\begin{metadefinition}[Teoria]\label{mdef:teoriaPrimeiraOrdem}\index{teoria}\index{axioma!não lógico}
    Uma \emph{teoria de primeira ordem} consiste em um léxico $\mathcal L$ para uma linguagem de primeira ordem munido de uma coleção de fórmulas $\Sigma$ de $\mathcal L$, chamadas de \emph{axiomas}.
    Tais axiomas também podem ser chamados de \emph{axiomas não lógicos}, para distingui-los dos axiomas lógicos enunciados na Metadefinição~\ref{mdef:axiomasLogicos}.
\end{metadefinition}

\begin{metadefinition}[Demonstração]\label{mdef:demonstracaoFormal}\index{demonstração formal}\index{teorema}\index{modus ponens}
    Seja $T$ uma teoria de primeira ordem e $\phi$ uma $\mathcal L$-fórmulas, em que $\mathcal L$ é o léxico de $T$.
    Uma \emph{demonstração} de $\phi$ a partir de $T$ é uma sequência finita de $\mathcal L$-fórmulas $\psi_0, \dots, \psi_n$ tal que para cada $i$:
    \begin{itemize}
        \item $\psi_i$ é um axioma lógico, ou
        \item $\psi_i$ é um axioma de $T$, ou
        \item $\psi_i$ é obtida por \emph{modus ponens} a partir de duas fórmulas anteriores.
        Ou seja, existem $j, k<i$ tais que $\psi_j$ é  $\psi_k \rightarrow \psi_i$.
    \end{itemize}

    Caso exista uma demonstração de $\phi$ a partir de $T$, dizemos que $\phi$ é um \emph{teorema} de $T$, e escrevemos $T \vdash \phi$.
    Caso os axiomas de $T$ sejam $\Sigma$ e o léxico de $T$ seja $\mathcal L$, também podemos escrever $\Sigma \vdash_{\mathcal L} \phi$ (e $\Sigma \nvdash_{\mathcal L} \phi$).
    Por fim, caso $\mathcal L$ esteja claro pelo contexto, podemos simplesmente escrever $\Sigma \vdash \phi$ (e $\Sigma \nvdash \phi$).
\end{metadefinition}

Uma \emph{regra de inferência} é uma regra formal que permite inferir como teorema novas fórmulas a partir de fórmulas já demonstradas.
No formalismo apresentado, a única regra de inferência é o \emph{modus ponens}, que permite inferir $\psi$ a partir de $\phi$ e $\phi \rightarrow \psi$.

Na prática, raramente escrevemos, em matemática, uma demonstração tão formal quanto a Metadefinição~\ref{mdef:demonstracaoFormal} exige.
Para reforçar a ideia de que as demonstrações usuais em matemática poderiam, em tese, ser reescritas de forma tão formal, apresentaremos alguns resultados relevantes.

Começaremos com alguns enunciados simples.
O primeiro resultado pode ser interpretado como a asserção de que resultados já provados podem ser usados para provar novos resultados.
\begin{metalemma}(Modus ponens)\label{mlem:modusPonens}
    Seja $T$ uma teoria de primeira ordem e $\phi, \psi$ $\mathcal L$-fórmulas, em que $\mathcal L$ é o léxico de $T$.
    Se $T \vdash \phi$ e $T \vdash \phi \rightarrow \psi$, então $T \vdash \psi$.
\end{metalemma}
\begin{proof}
    Existem demonstrações de $\phi$ e de $\phi \rightarrow \psi$ a partir de $T$, digamos $\psi_0, \dots, \psi_n$ e $\psi'_0, \dots, \psi'_m$, respectivamente, de modo que $\psi_n$ é $\phi$ e $\psi'_m$ é $\phi \rightarrow \psi$.
    Considere a seguinte sequência de fórmulas:
    \begin{equation*}
        \psi_0, \dots, \psi_n, \psi'_0, \dots, \psi'_{m}, \psi.
    \end{equation*}

    A sequência $\psi_0, \dots, \psi_n, \psi'_0, \dots, \psi'_{m}$, por ser uma concatenação de duas demonstrações, satisfaz claramente as condições para ser uma demonstração.
    Ao acrescentar $\psi$ ao final da sequência, ainda temos as condições satisfeitas, uma vez que $\psi$ é obtida a partir de \emph{modus ponens} de $\phi$ e $\phi \rightarrow \psi$, que são, respectivamente, $\psi_n$ e $\psi'_m$.
\end{proof}

O resultado a seguir nos diz que se uma teoria $T'$ prova todas as premissas de $T$, então $T'$ prova todos os teoremas de $T$.

\begin{metaproposition}\label{mprop:extensaoDeTeoria}
    Seja, $T, T'$ teorias de primeira ordem com léxicos $\mathcal L$ e $\mathcal L'$, respectivamente, tais que todo símbolo de $\mathcal L$ é um símbolo de $\mathcal L'$.
    Suponha que toda fórmula de $\Sigma$ é um teorema de $T'$.
    
    Então todo teorema de $T$ é um teorema de $T'$.
    Ou seja, para cada $\mathcal L$-fórmula $\phi$,

    \begin{equation*}
        \Sigma \vdash_\mathcal L \phi \implies \Sigma' \vdash_\mathcal L \phi.
    \end{equation*}

    Em particular, se $T'$ é uma teoria de primeira ordem obtida acrescentando novos axiomas a uma teoria $T$, então todo teorema de $T$ é um teorema de $T'$.
\end{metaproposition}

\begin{proof}
    Seja $\phi$ uma $\mathcal L$-fórmula tal que $\Sigma \vdash_\mathcal L \phi$.
    Então existe uma demonstração finita de $\phi$ a partir de $\Sigma$, $\phi_0, \dots, \phi_n$.
    Veremos, por indução em $m$, para $m\leq n$, que cada $\phi_m$ é um teorema de $T'$.

    Seja dado $m$, e suponha que o resultado é válido para cada $i<m$.
    
    \textbf{Caso 1:} $\phi_m$ é um axioma lógico.
    Nesse caso, $\phi_m$ é trivialmente um teorema de $T'$.

    \textbf{Caso 2:} $\phi_m$ é um axioma de $T$.
    Nesse caso, por hipótese, $\phi_m$ é um teorema de $T'$.

    \textbf{Caso 3:} $\phi_m$ é obtida a partir de \emph{modus ponens} de duas fórmulas anteriores, digamos $\phi_j$ e $\phi_k$, com $j, k<m$ e $\phi_j$ é $\phi_k \rightarrow \phi_m$.
    Por hipótese de indução, $T' \vdash \phi_k$ e $T' \vdash \phi_k\rightarrow \phi_m$.
    Do Metalema~\ref{mlem:modusPonens}, temos que $\phi_m$ é um teorema de $T'$.
\end{proof}

Façamos uma pausa para observar algo importante acontecendo na metateoria (isto é, no nosso raciocínio sobre demonstrações).
Ambos os resultados acima concluem que ``existem demonstrações''.
É importante observar que tais provas existenciais são construtivas: elas de fato fornecem um algoritmo para, dados os objetos declarados nas hipóteses, efetivamente redigir uma demonstração.

O próximo resultado nos mostra o motivo de contradições serem tão indesejáveis em Matemática: na Lógica Clássica, ao admitir uma, não estamos apenas sendo incoerentes, estamos abrindo margem para provar qualquer coisa: um verdadeiro desastre.

\begin{metadefinition}\label{mdef:inconsistencia}\index{inconsistência}
    Seja $T$ uma teoria de primeira ordem.
    Dizemos que $T$ é \emph{inconsistente} se existe uma fórmula $\phi$ tal que $T \vdash \phi$ e $T \vdash \neg \phi$.
\end{metadefinition}

\begin{metaproposition}[Princípio da explosão]\label{mprop:principioDaExplosao}\index{princípio da explosão}
    Seja $T$ uma teoria de primeira ordem.
    Se $T$ é inconsistente, então para cada fórmula $\psi$, $T \vdash \psi$.
\end{metaproposition}
\begin{proof}
    Tome $\psi$ uma fórmula arbitrária.
    Seja $\phi$ uma fórmula tal que $T \vdash \phi$ e $T \vdash \neg \phi$.

    Note que $A\rightarrow (\neg A\rightarrow B)$ é uma tautologia.
    Assim, $\phi \rightarrow (\neg \phi \rightarrow \psi)$ é uma instância de tautologia, e uma fórmula, e, portanto, um axioma lógico.
    Desta forma, $T \vdash \phi \rightarrow (\neg \phi \rightarrow \psi)$.

    Como $T \vdash \phi$, do Metalema~\ref{mlem:modusPonens} temos que $T \vdash \neg \phi \rightarrow \psi$.

    Como $T \vdash \neg \phi$, novamente do Metalema~\ref{mlem:modusPonens}, segue que $T \vdash \psi$.
\end{proof}

Abaixo, provaremos, como exemplo, um teorema da lógica pura que intuitivamente diz que o universo de discurso é não vazio.

\begin{metaproposition}\label{mprop:universoNaoVazio}
    Para cada teoria de primeira ordem $T$, $T \vdash \exists x (x=x)$.
\end{metaproposition}
\begin{proof}
    Considere a seguinte demonstração.

    \begin{align*}
    1.\ & \forall x (x=x) & \text{(axioma lógico 7.)} \\
    2.\ & \forall x (x=x \rightarrow \exists x (x=x)) & \text{(axioma lógico 5.)} \\
    3.\ & \forall x ((x=x) \rightarrow \exists x (x=x))\rightarrow (\forall x (x=x) \rightarrow \forall x \exists x (x=x)) & \text{(axioma lógico 3.)} \\
    4.\ & \forall x (x=x) \rightarrow \forall x \exists x (x=x) & \text{(modus ponens de 2. e 3.)} \\
    5.\ & \forall x\exists x (x=x) & \text{(modus ponens de 1. e 4.)} \\
    6.\ & \forall x\exists x (x=x) \rightarrow \exists x(x=x) & \text{(axioma lógico 2.)} \\
    7.\ & \exists x (x=x) & \text{(modus ponens de 5. e 6.)}
    \end{align*}
\end{proof}